\begingroup
\renewcommand{\arraystretch}{1.45}
\clearpage\fancyhead[L]{\small BEMERKUNGEN.}\pdfbookmark[0]{Bemerkungen zur Interpolation}{interpolation-bemerkungen}\SourceLeaf{337}
\begin{center}\textbf{BEMERKUNGEN.}\end{center}
\begin{center}\rule{3em}{0.3pt}\end{center}

Die vorliegende Abhandlung über Interpolation ist die erneute Ausarbeitung einer Untersuchung, von der ein früherer Entwurf im Oktober 1805 begonnen worden zu sein scheint. Sie enthält die Methoden, die bei der Berechnung der Ausdrücke für die durch Störungskräfte hervorgerufenen Änderungen der Elemente einer Planetenbahn vielfach angewandt worden sind. Dieser Ermittlung der Störungen steht ein anderes Verfahren zur Seite, das unmittelbar die Werte der Bahnelemente für einzelne Zeitpunkte ergibt und das, wie aus den im vorgefundenen handschriftlichen Nachlass aufgezeichneten abschließenden Formeln hervorgeht, auf folgenden Prinzipien beruht.

\Needspace{28\baselineskip}
Bilden die Werte des Arguments, für die die Werte der Funktion gegeben sind, eine arithmetische Folge, und hat man ein Argument $x$ eingeführt, das bei jenen Werten wie die ganzen Zahlen fortschreitet, bezeichnet die gegebenen Werte mit $fx$, führt deren aufeinanderfolgende Differenzenreihen ein:
\[
\begin{aligned}
&f(x+1)-fx=f^{1}(x+\frac12),\\[6pt]
&f^{1}(x+\frac12)-f^{1}(x-\frac12)=f^{2}x,\\[6pt]
&f^{2}(x+1)-f^{2}x=f^{3}(x+\frac12)\ \mbox{usw.}
\end{aligned}
\]
ferner deren mit beliebigen Anfangsgliedern beginnende Summenreihen
\hfill\break{}
$f^{-1}(x+\nobreak{}\tfrac12),\allowbreak{}\ f^{-2}x,\allowbreak{}\ f^{-3}(x+\nobreak{}\tfrac12)$ usw.,
für die
\[
\begin{aligned}
&fx=f^{-1}(x+\frac12)-f^{-1}(x-\frac12),\\[6pt]
&f^{-1}(x+\frac12)=f^{-2}(x+1)-f^{-2}x\ \mbox{usw.}
\end{aligned}
\]
schließlich die Mittelwerte der aufeinanderfolgenden Glieder der so entstandenen Reihen, nämlich
\[
\begin{aligned}
&\mathinner{.\,.\,.} f^{-1}x=\frac12f^{-1}(x+\frac12)+\frac12f^{-1}(x-\frac12),\\[6pt]
&f(x+\frac12)=\frac12f(x+1)+\frac12fx,\\[6pt]
&f^{1}x=\frac12f^{1}(x+\frac12)+\frac12f^{1}(x-\frac12)\mathinner{.\,.\,.}
\end{aligned}
\]
also in übersichtlicher Anordnung zusammengestellt das folgende System von gegebenen und daraus zunächst berechneten Werten
\SourceLeaf{338}

\begin{landscape}\thispagestyle{empty}\begingroup\fontsize{9}{12}\selectfont\setlength{\arraycolsep}{3pt}\renewcommand{\arraystretch}{1.3}
\[
\begin{array}{ccc|cc||cc||cc|ccc}
 & \mbox{M. W.} & \mbox{II. Summe.} & \mbox{M. W.} & \mbox{I. Summe} & \mbox{gegeb. W.} & \mbox{M. W.} & \mbox{I. Differ.} & \mbox{M. W.} & \mbox{II. Differ.} & \mbox{M. W.} & \\
 & & \cdot & \cdot & & \cdot & & & \cdot & \cdot & & \\
 & \cdot & & & \cdot & & \cdot & \cdot & & & \cdot & \cdot \\
 & & f^{-2}(x-1) & f^{-1}(x-1) & & f(x-1) & & & f^{1}(x-1) & f^{2}(x-1) & & \\
\cdot & f^{-2}(x-\tfrac12) & & & f^{-1}(x-\tfrac12) & & f(x-\tfrac12) & f^{1}(x-\tfrac12) & & & f^{2}(x-\tfrac12) & \cdot \\
 & & f^{-2}x & f^{-1}x & & fx & & & f^{1}x & f^{2}x & & \\
\cdot & f^{-2}(x+\tfrac12) & & & f^{-1}(x+\tfrac12) & & f(x+\tfrac12) & f^{1}(x+\tfrac12) & & & f^{2}(x+\tfrac12) & \cdot \\
 & & f^{-2}(x+1) & f^{-1}(x+1) & & f(x+1) & & & f^{1}(x+1) & f^{2}(x+1) & & \\
\cdot & \cdot & & & \cdot & & \cdot & \cdot & & & \cdot & \cdot \\
 & & \cdot & \cdot & & \cdot & & & \cdot & \cdot & &
\end{array}
\]
\endgroup\end{landscape}

\Needspace{30\baselineskip}
so gibt Artikel 4, wenn man $1.$ für $t,\allowbreak{}\ a,\allowbreak{}\ b,\allowbreak{}\ c,\allowbreak{}\ d,\allowbreak{}\ e\ldotp\ldotp$
der Reihe nach $x+t,\allowbreak{}\ x,\allowbreak{}\ x+1,\allowbreak{}\ x-1,\allowbreak{}\ x+2,\allowbreak{}\ x-2$, $2.$ für
$t,\allowbreak{}\ a,\allowbreak{}\ b,\allowbreak{}\ c,\allowbreak{}\ d,\allowbreak{}\ e\ldotp\ldotp$ der Reihe nach $x+t,\allowbreak{}\ x,\allowbreak{}\ x-1,\allowbreak{}\ x+1,\allowbreak{}\ x-2,\allowbreak{}\ x+2\mathinner{.\,\allowbreak{}.\,\allowbreak{}.}$ einsetzt,
$3.$ aus den entsprechenden Seiten der beiden so erhaltenen Gleichungen die halbe Summe bildet,
$4.$ aus denselben Gleichungen, nachdem zuvor in der zweiten $x+1$ und $t-1$ der Reihe nach für $x$ und $t$ eingesetzt worden sind,
die halbe Summe bildet und $5.$ in der zuletzt erhaltenen Gleichung $t=\tfrac12$ setzt,
die folgenden Gleichungen für die durch Interpolation zu bestimmenden Werte $\varphi(x+\nobreak{}t)$ und
$\varphi(x+\nobreak{}\tfrac12)$ derjenigen Funktion $\varphi x$, die bei jedem ganzzahligen Wert
von $x$ der gegebenen Größe $fx$ gleich wird:
\[
\begin{aligned}
&\varphi(x+t)={}fx+t\cdot f^{1}(x+\frac12)+\frac{t\cdot t-1}{1\cdot2}f^{2}x
\\[6pt]
&\qquad{}+\frac{t\cdot t-1\cdot t+1}{1\cdot2\cdot3}f^{3}(x+\frac12)+\ldotp\ldotp\\[6pt]
&+\frac{\Pi(t+n-1)}{\Pi 2n\cdot\Pi(t-n-1)}f^{2n}x
\\[6pt]
&\qquad{}+\frac{\Pi(t+n)}{\Pi(2n+1)\cdot\Pi(t-n-1)}f^{2n+1}(x+\frac12)\ldotp\ldotp
\end{aligned}
\]
\Needspace{18\baselineskip}
\[
\begin{aligned}
&\varphi(x+t)={}fx+t\cdot f^{1}(x-\frac12)+\frac{t\cdot t+1}{1\cdot2}f^{2}x
\\[6pt]
&\qquad{}+\frac{t\cdot t+1\cdot t-1}{1\cdot2\cdot3}f^{3}(x-\frac12)+\ldotp\ldotp\\[6pt]
&+\frac{\Pi(t+n)}{\Pi 2n\cdot\Pi(t-n)}f^{2n}x
\\[6pt]
&\qquad{}+\frac{\Pi(t+n)}{\Pi(2n+1)\cdot\Pi(t-n-1)}f^{2n+1}(x-\frac12)\mathinner{.\,.\,.}
\end{aligned}
\]
\[
\begin{aligned}
&\varphi(x+t)={}fx+t\cdot f^{1}x+\frac{t\cdot t}{1\cdot2}f^{2}x
+\frac{t\cdot t-1\cdot t+1}{1\cdot2\cdot3}f^{3}x+\ldotp\ldotp\\[6pt]
&+\frac{\Pi(t+n-1)}{\Pi 2n\cdot\Pi(t-n)}t f^{2n}x
+\frac{\Pi(t+n)}{\Pi(2n+1)\Pi(t-n-1)}f^{2n+1}x\\[6pt]
&\qquad{}+\mathinner{.\,.\,.}
\end{aligned}
\]
\Needspace{22\baselineskip}
\[
\begin{aligned}
&\varphi(x+t)={}f(x+\frac12)+(t-\frac12)f^{1}(x+\frac12)
+\frac{t\cdot t-1}{1\cdot2}f^{2}(x+\frac12)
\\[6pt]
&\qquad{}+\frac{t\cdot t-1}{1\cdot2}(t-\frac12)f^{3}(x+\frac12)+\ldotp\ldotp\\[6pt]
&+\frac{\Pi(t+n-1)}{\Pi 2n\cdot\Pi(t-n-1)}f^{2n}(x+\frac12)
\\[6pt]
&\qquad{}+\frac{\Pi(t+n-1)}{\Pi(2n+1)\cdot\Pi(t-n-1)}(t-\frac12)f^{2n+1}(x+\frac12)+\ldotp\ldotp
\end{aligned}
\]
\[
\begin{aligned}
&\varphi(x+\frac12)={}f(x+\frac12)-\frac18f^{2}(x+\frac12)
+\frac{3}{128}f^{4}(x+\frac12)\\[6pt]
&\qquad{}-\frac{5}{1024}f^{6}(x+\frac12)
+\frac{35}{32768}f^{8}(x+\frac12)\mathinner{.\,.\,.}\\[6pt]
&+(-1)^n\frac{\Pi 2n}{\Pi n\cdot\Pi n}\,2^{-4n}\cdot f^{2n}(x+\frac12)+\mathinner{.\,.\,.}
\end{aligned}
\]
\Needspace{14\baselineskip}
Die Koeffizienten von $f^m$ in der dritten, vierten und fünften Gleichung sind gleich den Koeffizienten von
$(2i\sin\omega)^m$ in den Reihenentwicklungen der Reihe nach für
$\cos 2t\omega+i\dfrac{\sin 2t\omega}{\cos\omega},\allowbreak{}\allowbreak \quad
\dfrac{\cos(2t-\nobreak{}1)\omega}{\cos\omega}+i\sin(2t-\nobreak{}1)\omega$
und $\dfrac{1}{\cos\omega}$ nach Potenzen der Größe $2i\sin\omega$.
\Needspace{28\baselineskip}
\par\SourceLeaf{339}
Setzt man nun
\[
\begin{aligned}
&\Phi(x+\frac12)={}f^{-1}(x+\frac12)+\frac{1}{24}f^{1}(x+\frac12)
-\frac{17}{5760}f^{3}(x+\frac12)\\[6pt]
&\qquad{}+\frac{367}{967680}f^{5}(x+\frac12)
-\frac{27859}{164486400}f^{7}(x+\frac12)\\[6pt]
&+\frac{1295803}{122624409600}f^{9}(x+\frac12)-\ldotp\ldotp
\end{aligned}
\]
\[
\begin{aligned}
&\Phi x={}f^{-1}x-\frac{1}{12}f^{1}x+\frac{11}{720}f^{3}x
-\frac{191}{60480}f^{5}x+\frac{2497}{3628800}f^{7}x\\[6pt]
&-\frac{14797}{95800320}f^{9}x+\frac{92427157}{2615348736000}f^{11}x-\mathinner{.\,.\,.}
\end{aligned}
\]
\[
\begin{aligned}
&\Psi x={}f^{-2}x+\frac{1}{12}fx-\frac{1}{240}f^{2}x
+\frac{31}{60480}f^{4}x-\frac{289}{3628800}f^{6}x\\[6pt]
&+\frac{317}{22809600}f^{8}x-\ldotp\ldotp
\end{aligned}
\]
\[
\begin{aligned}
&\Psi(x+\frac12)=f^{-2}(x+\frac12)-\frac{1}{24}f(x+\frac12)
+\frac{17}{1920}f^{2}(x+\frac12)\\[6pt]
&\qquad{}-\frac{367}{193536}f^{4}(x+\frac12)+\ldotp-\ldotp\ldotp
\end{aligned}
\]
\Needspace{30\baselineskip}
indem man als Koeffizienten von $f^m$ in diesen vier Gleichungen die Koeffizienten nimmt, die bei den
Entwicklungen der Reihe nach für
$\dfrac{1}{2i\omega},\allowbreak{}\allowbreak \ \dfrac{1}{2i\omega\cos\omega},\allowbreak{}\allowbreak \ -\dfrac{1}{4\omega\omega},\allowbreak{}\allowbreak \ -\dfrac{1}{4\omega\omega\cos\omega}$
noch Potenzen von $2i\sin\omega$ entstehen,
so sind $\Phi(x+\nobreak{}\tfrac12),\allowbreak{}\ \Phi x,\allowbreak{}\ \Psi x,\allowbreak{}\ \Psi(x+\nobreak{}\tfrac12)$ die besonderen Werte solcher Funktionen
$\Phi(x+\nobreak{}t),\allowbreak{}\ \Psi(x+\nobreak{}t)$,
für die die Gleichungen
\[
\begin{aligned}
&\Phi(x+t)-\Phi(x+t_0)=\int_{t_0}^{t}\varphi(x+t)\,dt),\\[6pt]
&\Psi(x+t)-\Psi(x+t_1)=\int_{t_1}^{t}\Phi(x+t)\,dt
\end{aligned}
\]
gelten; sie werden also bestimmten Integralen von $\varphi(x+\nobreak{}t)\,\allowbreak{}dt$ und
$\varphi(x+\nobreak{}t)\cdot dt\cdot dt$ gleich, wenn
man die Anfangswerte der Summenreihen
\hfill\break{}
$f^{-1}(x-\nobreak{}\tfrac12),\allowbreak{}\ f^{-1}(x+\nobreak{}\tfrac12)\ldotp\ldotp,\allowbreak{}\ f^{-2}x,\allowbreak{}\ f^{-2}(x+\nobreak{}1)\ldotp\ldotp$
auf geeignete Weise auswählt.

Die Herleitung der Reihen für $\Phi$ und $\Psi$ aus der dritten und vierten der obigen Gleichungen für
$\varphi(x+\nobreak{}t)$ erhält man unmittelbar, wenn man die Integrationen so ausführt, dass zunächst die Ausdrücke für
\hfill\break{}
$\Phi(x+\nobreak{}\tfrac12)-\Phi(x-\nobreak{}\tfrac12),\allowbreak{}\quad \Phi(x+\nobreak{}1)-\Phi x,\allowbreak{}\quad
\Psi(x+\nobreak{}1)-2\Psi x+\Psi(x-\nobreak{}1)$ und
$\Psi(x+\nobreak{}\tfrac32)-2\Psi(x+\nobreak{}\tfrac12)+\Psi(x-\nobreak{}\tfrac12)$
entstehen.
\begin{flushright}\textsc{Schering.}\end{flushright}
\begin{center}\rule{9em}{0.3pt}\end{center}
\clearpage
\endgroup
